\documentclass[11pt]{article}
\usepackage[margin=1in]{geometry}
\usepackage{amsmath,amssymb,booktabs,graphicx,hyperref,xcolor}
\usepackage[numbers,sort&compress]{natbib}
\graphicspath{{figures/}}
\newcommand{\fig}[1]{\includegraphics[width=\linewidth]{#1.pdf}}

\title{Resource-Aware Evaluation of Quantum Error Mitigation Techniques\\Under Noisy Circuit Execution}
\author{Nagarjun Gopalaswamy\\Independent researcher}
\date{Draft manuscript --- not submitted, not peer reviewed}

\begin{document}
\maketitle

\begin{abstract}
Quantum error mitigation (QEM) improves expectation-value estimates from noisy quantum circuits, but it spends extra
calibration circuits, noise-amplified executions and shots. We ask how readout error mitigation (REM), zero-noise
extrapolation (ZNE; linear and Richardson) and their combination trade estimation accuracy against statistical uncertainty
and sampling cost, as circuit depth, noise intensity and measurement budget vary. We built QEM-Bench, an open,
reproducible harness whose resource ledger counts every executed shot. We evaluated Bell, GHZ (3--8 qubits) and small QAOA
Max-Cut circuits under phenomenological Qiskit Aer noise at shot budgets from 256 to 16{,}384, with 30 independent
repetitions per condition and two accounting rules: equal primary shots (A) and equal total budget including calibration (B).
In 8{,}532 simulated conditions (255{,}960 estimates), REM+ZNE-linear had the lowest median RMSE: $-81\%$ relative to
baseline under A, where it used $4.3\times$ the shots, and $-68\%$ at equal total budget. REM alone gave $-60\%$ ($-54\%$ under B) for
$2.3\times$ the shots under A. ZNE alone gave about $-22\%$ ($-19\%$ under B). Under the established shot-normalised improvement
factor, ZNE alone beat simply spending the extra shots on the baseline in only 10--18\% of conditions. At the smallest budgets with low noise,
3\% of equal-budget conditions showed no significant benefit from any method, and Richardson ZNE was significantly worse than baseline in 17\% of
equal-budget conditions at 256 shots. We also found that unweighted partial gate folding can mis-scale noise severely on small
circuits. All results are simulations with simplified noise, not hardware measurements.
\end{abstract}

\section{Introduction}
Noisy intermediate-scale devices suffer gate, decoherence and measurement errors. QEM techniques reduce the bias of
estimated observables without fault tolerance~\cite{temme2017,li2017,endo2018,cai2023review}, at a sampling cost that is
unavoidable in general~\cite{takagi2022,quek2024}. Accuracy alone therefore cannot rank techniques: a method that halves the error
but needs four times the shots may lose to the unmitigated estimator given those same shots. Resource-normalised comparison has
precedent~\cite{russo2023}. We extend it to readout mitigation and REM+ZNE combinations, using two explicit accounting rules,
paired statistics with false-discovery-rate control and Pareto analysis.

\paragraph{Research questions.} RQ1: REM benefit versus readout error. RQ2: ZNE versus depth and gate noise. RQ3: whether
REM+ZNE helps consistently. RQ4: behaviour under constrained shot budgets. RQ5: regions where mitigation costs resources without meaningful
gain. RQ6: resource-aware ranking.

\paragraph{Contributions (scoped to this simulated study).} (i) An open benchmarking harness with an enforceable
resource-accounting validator. (ii) A controlled dataset and analysis of four mitigation configurations under equal-shot and
equal-budget rules. (iii) A dimensionless efficiency metric $\eta$, placed alongside the established improvement factor $\mu$~\cite{russo2023}.
(iv) A documented noise-scaling pitfall, with its correction. We do not claim to be the first comparison of these techniques.

\section{Related Work}
ZNE was introduced through Richardson extrapolation of noise-amplified expectation values~\cite{temme2017,li2017} and
demonstrated on hardware using pulse stretching~\cite{kandala2019}. Digital unitary folding~\cite{giurgica2020} and software frameworks such as
Mitiq~\cite{larose2022} made ZNE available at the gate level. Exponential extrapolation and PEC with learned noise were
developed in~\cite{endo2018}, and PEC with sparse Pauli--Lindblad models scales to large processors~\cite{vandenberg2023},
which also supports large-scale ZNE~\cite{kim2023}. Readout mitigation by detector tomography and matrix inversion~\cite{maciejewski2020} has
scalable variants~\cite{bravyi2021,nation2021}. Learning-based mitigation includes Clifford data regression~\cite{czarnik2021}.
Fundamental sampling-cost bounds are given in~\cite{takagi2022,quek2024}, and optimised Richardson node choice and shot allocation in~\cite{krebsbach2022}.
The closest prior work, Russo et al.~\cite{russo2023}, benchmarks ZNE and PEC across platforms using a shot-normalised
improvement factor. Reviews~\cite{endo2021,cai2023review} discuss combining techniques and the bias--variance trade-off qualitatively.

\section{Background}
For an observable $O$ diagonal after basis rotation, the baseline estimator is $\hat\theta=\sum_x\hat p(x)f(x)$, with variance
$\mathrm{Var}_p[f]/N$. \emph{Readout error} acts as a column-stochastic assignment matrix $A$, giving $p_{\rm meas}=Ap$. REM
estimates $A$ from basis-state calibration circuits and evaluates $\hat\theta=(A^{-T}f)\cdot\hat p$. This estimator is unbiased for the linear functional
when $A$ is known, but its variance is inflated. \emph{ZNE} measures $E(\lambda)$ at scale factors $\lambda_i$, obtained here by folding
$G\to G(G^\dagger G)^k$, and extrapolates $E_0=\sum_i\gamma_iE(\lambda_i)$. Linear least squares and Richardson interpolation are both linear in
the data, so $\mathrm{Var}(E_0)=\sum_i\gamma_i^2\mathrm{Var}(E_i)$. For $\lambda=(1,3,5)$, Richardson gives
$\gamma=(15/8,-5/4,3/8)$, $\sum\gamma_i^2\approx4.2$. The \emph{sampling overhead} is $O=N^{\rm tot}_{\rm mit}/N^{\rm tot}_{\rm base}$.

\section{Methodology}
\begin{figure}[t]\fig{fig01_architecture}\caption{QEM-Bench pipeline.}\label{fig:arch}\end{figure}
\textbf{Circuits.} Bell ($|\Phi^+\rangle$ via CX or CZ, and a depth-padded variant with two barrier-protected CX--CX pairs), with observables
$XX, ZZ, YY$. GHZ with $n\in\{3,4,5,6,8\}$, observables $X^{\otimes n}$ and $Z_0Z_{n-1}$. QAOA Max-Cut on seeded connected $G(n,0.5)$ graphs, $n\in\{4,6\}$,
$p\in\{1,2\}$, with angles optimised noiselessly and then fixed, and the normalised cut $\langle C\rangle/|E|$ as observable. Exact references come from statevector
simulation, and Max-Cut is solved by brute force. Circuits are transpiled to $\{rz,sx,x,cx\}$ (all-to-all connectivity).

\textbf{Noise.} Phenomenological depolarizing noise on $sx,x$ ($\lambda_1$) and $cx$ ($\lambda_2$), plus independent asymmetric
readout error ($p_{1|0}=0.8r$, $p_{0|1}=1.2r$). The levels $(\lambda_1,\lambda_2,r)$ are low $(5\!\times\!10^{-4},5\!\times\!10^{-3},5\!\times\!10^{-3})$,
moderate $(2\!\times\!10^{-3},0.015,0.02)$ and high $(5\!\times\!10^{-3},0.03,0.05)$. Sensitivity sweeps scale gate noise ($0$--$6\times$) or readout
noise ($0$--$6\times$) around the moderate level. Variants use symmetric readout or add T1/T2 thermal relaxation. These settings are
experimental design choices, not hardware claims.

\textbf{Mitigation.} REM uses full $2^n$ calibration for $n\le3$ and a tensored two-circuit model otherwise. It applies an unconstrained inverse,
records the condition number, and fails explicitly above $10^{12}$. ZNE uses global folding with $\lambda=(1,3,5)$ in the core grid, and local folding,
$(1,2,3)$ and $(1,3,5,7)$ in a design block. Partial folds are chosen by error-weighted gate selection, and extrapolation uses the
\emph{measured} error-mass scale $\lambda_{\rm eff}$. REM+ZNE shares one calibration and corrects every scale point before extrapolating, since
measurements are not folded.

\textbf{Resource accounting.} With $N$ shots per circuit, $S$ settings, $K$ scales and $M$ calibration circuits: under comparison A the baseline spends
$NS$, REM $NS+MN$, ZNE $NSK$ and REM+ZNE $NSK+MN$. Under comparison B every method spends at most $T=NS$. REM gives 20\% of $T$ to calibration,
and ZNE allocates shots in proportion to $|\gamma_i|$~\cite{krebsbach2022}. A validator checks these identities on every dataset row.

\textbf{Metrics and statistics.} Bias, variance, MSE, RMSE and MAE, with 95\% percentile-bootstrap CIs~\cite{efron1993} (2000 resamples) and 30
repetitions per condition, using independent \texttt{SeedSequence} streams. Methods are compared by two-sided Wilcoxon signed-rank tests on
paired absolute errors (common random numbers under A; paired by repetition under B), with the rank-biserial effect size, paired-bootstrap
CIs of $\Delta$MSE and Benjamini--Hochberg control~\cite{bh1995} over all 12{,}798 tests (Holm~\cite{holm1979} also reported). We say ``practical''
for $|\Delta\mathrm{RMSE}|\ge10\%$. The efficiency metric is
\begin{equation}
\eta=\frac{(E_b-E_m)/E_b}{(C_m-C_b)/C_b},
\end{equation}
with $E$ the RMSE and $C$ the total shots. It is dimensionless, and edge cases ($\Delta C\le0$, $E_b=0$) are labelled explicitly. We also report $\mu=(E_b/E_m)/\sqrt{C_m/C_b}$~\cite{russo2023}.
Note that $\mu>1$ is algebraically equivalent to a work-normalised MSE ratio below 1. The pre-specified decision rules for H1--H5 are in the supplement.

\section{Experimental Setup}
Python 3.11.18, Qiskit 2.2.3, Qiskit Aer 0.17.2, NumPy 2.4.6, SciPy 1.17.1, pandas 3.1.1, running on a 4-core aarch64 Linux container.
Root seed 20260601, transpiler optimisation level 1. Sampling in the standard run is \emph{exact-distribution}: Aer density-matrix simulation plus the exact readout channel,
followed by multinomial sampling. A pilot (3 instances, 2 noise levels, $N\in\{1024,4096\}$, 30 repetitions) compared this mode with
per-shot Aer sampling over 288 conditions. No mean difference was significant after FDR correction (4.9\% raw $p<0.05$, consistent with the nominal 5\%), and
1 of 288 variance tests was significant. The pilot also provided the wall-clock timing study (Table~\ref{tab:timing}). The executed
\texttt{standard} profile is a reduced factorial: a core grid of 12 instances $\times$ 3 levels $\times$ 7 budgets, two sweeps, noise variants and three ZNE-design
blocks. The \texttt{full} profile is documented but was not executed.

\section{Results}
\begin{figure}[t]\fig{fig04_method_comparison}\caption{RMSE per instance, moderate noise, $N=4096$. Bars show 95\% bootstrap CIs.}\label{fig:methods}\end{figure}
\begin{table}[t]
\centering
\small
\caption{Core grid: median relative RMSE change versus baseline, fraction of conditions significantly better/worse (BH $q<0.05$) and median sampling overhead.}
\label{tab:core}
\resizebox{\linewidth}{!}{%
\begin{tabular}{llllll}
\toprule
comparison & method & median\_rel\_rmse\_change & frac\_better & frac\_worse & sampling\_overhead \\
\midrule
A & combined\_linear & -0.806 & 0.99 & 0 & 4.33 \\
A & combined\_richardson & -0.738 & 0.938 & 0.00621 & 4.33 \\
A & readout & -0.597 & 1 & 0 & 2.33 \\
A & zne\_linear & -0.222 & 0.99 & 0 & 3 \\
A & zne\_richardson & -0.223 & 0.905 & 0.0145 & 3 \\
B & combined\_linear & -0.681 & 0.932 & 0.00207 & 0.999 \\
B & combined\_richardson & -0.553 & 0.797 & 0.0393 & 1 \\
B & readout & -0.536 & 0.957 & 0 & 1 \\
B & zne\_linear & -0.195 & 0.876 & 0.00207 & 1 \\
B & zne\_richardson & -0.188 & 0.739 & 0.0393 & 1 \\
\bottomrule
\end{tabular}}
\end{table}


\textbf{Overall (RQ3, RQ6).} Table~\ref{tab:core} and Fig.~\ref{fig:methods} summarise the core grid. REM+ZNE-linear had the best mean
RMSE rank under both A (1.33 of 6) and B (1.58). REM ranked second under B. Under A, REM+ZNE was significantly better than
\emph{both} of its components in 69\% of conditions; under B, where its budget is split three ways, this fell to 40\% (Fig.~\ref{fig:comb}).

\textbf{Readout mitigation (RQ1, H1).} In the readout sweep, REM's median RMSE reduction grew monotonically from 7\% (no
readout error) to 89\% ($6\times$ moderate readout error), and it was significantly better than baseline with lower $|$bias$|$ in all 40
readout-dominated conditions (H1 SUPPORTED). The 7\% gain with \emph{zero} readout error is an artefact: noisy $x$ gates in the
calibration circuits are attributed to readout, and the partial inversion happens to counteract depolarizing contraction. The REM
variance was 1.23$\times$ the shot-noise prediction (median), because calibration uncertainty is not included in the analytic error bar.

\begin{figure}[t]\fig{fig03_noise_vs_error}\caption{Sensitivity sweeps over gate-noise (top) and readout-noise (bottom) multipliers (comparison A, $N=4096$).}\label{fig:sweeps}\end{figure}
\textbf{ZNE (RQ2, H2).} Under A, ZNE was significantly better than baseline in 100\% of moderate-noise conditions with $N\ge4096$, but in only
78\% at $N=256$. Contrary to the noise component of H2, the fraction \emph{rose} with noise level (low 87\%, moderate 98\%, high 99\%), and in the
gate sweep the relative improvement of Richardson ZNE grew with the gate-noise multiplier for Bell and GHZ-4 (up to $6\times$), saturating at 3--4$\times$ for GHZ-6 and QAOA. In this regime there is more bias to remove,
and the folded decay stays well resolved. H2 is supported by the pre-specified rule only through its budget component. ZNE left a substantial residual bias
(median $|$bias$|$ 0.08 versus 0.12 for baseline, moderate noise, $N=16384$), consistent with near-exponential depolarizing decay. Combined with REM,
the residual bias fell to 0.002--0.005.

\begin{figure}[t]\fig{fig06_shots_vs_rmse}\caption{Shot budget versus RMSE. Top: comparison A. Bottom: comparison B (equal total budget).}\label{fig:shots}\end{figure}

\textbf{Shot budgets (RQ4) and resource normalisation (RQ5, H5).} Fig.~\ref{fig:shots} shows that baseline RMSE plateaus at its bias while mitigated RMSE keeps falling
with $N$. At equal total budget and $N=256$, Richardson-based methods had \emph{higher} median RMSE than baseline (ZNE-Richardson $+5\%$,
REM+ZNE-Richardson $+21\%$), and ZNE-Richardson was significantly worse in 17\% of those conditions. Across the core grid, 3.3\% of equal-budget
conditions (all low noise, $N\le1024$) showed no significant benefit from any method (H5 SUPPORTED in its weak form). No condition had every method
significantly worse. Under the shot-normalised factor $\mu$, ZNE-linear exceeded 1 in only 10\% of comparison-A conditions (Richardson 18\%), versus 84\% for REM and
89\% for REM+ZNE-linear. In other words, by this criterion ZNE alone was usually less efficient than spending the extra shots on the
baseline, even though its accuracy gain was statistically significant. The median $\eta$ was 0.46 for REM, 0.24 for REM+ZNE-linear and 0.11 for ZNE (Fig.~\ref{fig:eff}).
Baseline points made up 16\% of the core Pareto frontier, all at the smallest shot totals; REM+ZNE-linear made up 48\% (Fig.~\ref{fig:pareto}).

\begin{figure}[t]\fig{fig08_efficiency_heatmap}\caption{Efficiency $\eta$ (comparison A, clipped) and RMSE reduction at equal budget (comparison B).}\label{fig:eff}\end{figure}
\begin{figure}[t]\fig{fig09_pareto_frontier}\caption{Pareto frontiers (RMSE versus total shots) across methods and budgets.}\label{fig:pareto}\end{figure}

\textbf{Depth (H4).} The pre-specified pooled test found no positive correlation between two-qubit gate count and baseline $|$bias$|$
(Spearman $\rho=-0.25$ to $-0.42$, $p>0.16$; H4 NOT SUPPORTED). The test design pooled observables with different bias scales: QAOA's
normalised cut tends towards 0.5 under depolarizing noise, so its absolute bias is small despite deep circuits. An exploratory within-family
analysis found perfectly monotone growth for GHZ ($\rho=1$ at every level), and state fidelity fell from 0.99 (Bell) to 0.61 (QAOA $n=6,p=2$) at moderate
noise (Fig.~\ref{fig:depth}). Because this analysis is post hoc, it does not change the verdict.

\begin{figure}[t]\fig{fig02_depth_vs_error}\caption{Two-qubit gate count versus RMSE (comparison A, $N=16384$).}\label{fig:depth}\end{figure}
\begin{figure}[t]\fig{fig12_combined_vs_individual}\caption{Combined mitigation versus the better individual method.}\label{fig:comb}\end{figure}

\textbf{ZNE design and noise scaling.} With unweighted gate-count partial folding (as in common digital-ZNE practice), our validation
measured $\lambda_{\rm eff}=1.18$ and $2.93$ for nominal $\lambda=2$ on GHZ-4, because the virtual $rz$ gates and the $sx$/$cx$ imbalance distort the count.
Error-weighted selection reduced the range to $[2.00,2.05]$ across the design block, and extrapolation uses the measured values (a regression test guards this).
Odd-integer folding scaled the error mass exactly (maximum deviation $2\times10^{-14}$) and preserved the unitary in every case. Median ZNE gains were similar across
designs (Fig.~\ref{fig:zne}). Adding the node $\lambda=7$ reduced the linear-fit gain (16\% versus 21\%), consistent with curvature in the decay.

\begin{figure}[t]\fig{fig11_zne_scale_factor}\caption{Expectation values versus scale factor, and ZNE design comparison.}\label{fig:zne}\end{figure}

\textbf{Families and noise variants.} \begin{table}[t]
\centering
\small
\caption{Median relative RMSE change vs baseline by circuit family (N=4096).}
\label{tab:family}
\resizebox{\linewidth}{!}{%
\begin{tabular}{lllllll}
\toprule
family & comparison & combined\_linear & combined\_richardson & readout & zne\_linear & zne\_richardson \\
\midrule
bell & A & -0.868 & -0.808 & -0.754 & -0.202 & -0.17 \\
bell & B & -0.784 & -0.676 & -0.664 & -0.187 & -0.18 \\
ghz & A & -0.888 & -0.851 & -0.678 & -0.23 & -0.27 \\
ghz & B & -0.809 & -0.719 & -0.646 & -0.237 & -0.233 \\
qaoa & A & -0.712 & -0.809 & -0.315 & -0.371 & -0.44 \\
qaoa & B & -0.663 & -0.671 & -0.302 & -0.317 & -0.431 \\
\bottomrule
\end{tabular}}
\end{table}
 REM helped QAOA least (median $-32\%$ versus $-75\%$ for Bell under A), because the normalised cut is
less sensitive to readout flips. ZNE helped QAOA most. Results under symmetric readout and T1/T2 relaxation were qualitatively unchanged (supplement).
\textbf{Timing.} In per-shot Aer sampling, mitigation multiplied wall-clock time by about 2.6 (REM), 3.6 (ZNE) and 5.2 (REM+ZNE) relative to baseline
(median, Table~\ref{tab:timing}). The classical post-processing time was below 1 ms. These timings are machine-specific.
\begin{table}[t]
\centering
\small
\caption{Wall-clock time per condition-repetition (s), per-shot Aer sampling, N=4096, comparison A (machine-specific).}
\label{tab:timing}
\begin{tabular}{lllllll}
\toprule
instance & method & t\_total\_mean & t\_total\_sd & t\_sampling\_mean & t\_calibration\_mean & t\_postprocess\_mean \\
\midrule
bell\_cx & baseline & 0.0211 & 0.0067 & 0.0210 & 0.0000 & 0.0000 \\
bell\_cx & combined\_linear & 0.0898 & 0.0046 & 0.0643 & 0.0252 & 0.0003 \\
bell\_cx & combined\_richardson & 0.0897 & 0.0046 & 0.0643 & 0.0252 & 0.0002 \\
bell\_cx & readout & 0.0463 & 0.0069 & 0.0210 & 0.0252 & 0.0000 \\
bell\_cx & zne\_linear & 0.0646 & 0.0072 & 0.0643 & 0.0000 & 0.0003 \\
bell\_cx & zne\_richardson & 0.0645 & 0.0072 & 0.0643 & 0.0000 & 0.0002 \\
ghz\_4 & baseline & 0.0173 & 0.0045 & 0.0172 & 0.0000 & 0.0000 \\
ghz\_4 & combined\_linear & 0.0719 & 0.0082 & 0.0536 & 0.0177 & 0.0006 \\
ghz\_4 & combined\_richardson & 0.0720 & 0.0081 & 0.0536 & 0.0177 & 0.0007 \\
ghz\_4 & readout & 0.0352 & 0.0051 & 0.0172 & 0.0177 & 0.0002 \\
ghz\_4 & zne\_linear & 0.0538 & 0.0087 & 0.0536 & 0.0000 & 0.0002 \\
ghz\_4 & zne\_richardson & 0.0537 & 0.0087 & 0.0536 & 0.0000 & 0.0001 \\
qaoa\_n4\_p1 & baseline & 0.0089 & 0.0019 & 0.0089 & 0.0000 & 0.0000 \\
qaoa\_n4\_p1 & combined\_linear & 0.0508 & 0.0052 & 0.0338 & 0.0166 & 0.0004 \\
qaoa\_n4\_p1 & combined\_richardson & 0.0507 & 0.0052 & 0.0338 & 0.0166 & 0.0003 \\
qaoa\_n4\_p1 & readout & 0.0256 & 0.0037 & 0.0089 & 0.0166 & 0.0001 \\
qaoa\_n4\_p1 & zne\_linear & 0.0339 & 0.0045 & 0.0338 & 0.0000 & 0.0001 \\
qaoa\_n4\_p1 & zne\_richardson & 0.0339 & 0.0045 & 0.0338 & 0.0000 & 0.0001 \\
\bottomrule
\end{tabular}
\end{table}


\section{Discussion}
In this simulated regime, the readout error contributed a large and easily removable share of the bias, so REM gave the best
efficiency per extra shot, and REM+ZNE the best accuracy. ZNE removed part of the gate-induced bias, but its variance amplification
($\sum\gamma_i^2\approx4.2$ for Richardson) and residual extrapolation bias meant that, after normalising for shots, it seldom paid off on its own.
This is consistent with the theoretical cost of mitigation~\cite{takagi2022,cai2023review}, and with the platform-dependent ZNE results of~\cite{russo2023}.
The accounting rule changed the conclusions: comparison A overstates the benefit of ZNE and of the combination relative to an
equal-budget view, and the most budget-sensitive configurations were the Richardson-based ones. Our efficiency metric $\eta$ and the established $\mu$
rank methods similarly, but they answer different questions. $\eta$ measures relative accuracy gained per relative extra cost. $\mu$ asks whether mitigation beats
spending the same shots on the baseline, and it implicitly assumes that baseline error falls as $1/\sqrt N$, which fails when the baseline is bias-dominated.

\section{Threats to Validity}
(1) \emph{Simulator and noise model:} depolarizing plus uncorrelated readout noise is favourable to tensored REM and to folding-based
ZNE (gate noise scales exactly with folding). Coherent, correlated, non-Markovian and drifting noise is absent. (2) \emph{Scale:} at most 8 qubits and
no routing, far from the asymptotic regime of~\cite{quek2024}. (3) \emph{Statistics:} 30 repetitions per condition, and
bootstrap CIs can be slightly anti-conservative (empirical coverage lay within the 0.90--0.98 acceptance band of our synthetic test). Conditions share instances, so tests are not
independent, although BH remains valid under positive dependence. (4) \emph{Extrapolation choices:} the scale sets, linear and Richardson
models only (no exponential fit), and a fixed calibration fraction of 20\%. (5) \emph{Accounting definitions:} conclusions under comparison B depend on
the calibration fraction and allocation rule. (6) \emph{Hypothesis rules} were fixed by the author without external review, and the H4 rule was poorly
designed (pooled observables). (7) \emph{Generalisation to hardware} is not established: no hardware experiments were run.

\section{Conclusion}
Under controlled simulated noise, REM and REM+ZNE reduced estimation error substantially and robustly, even at equal total budget,
whereas ZNE alone yielded statistically significant but resource-inefficient gains, and Richardson variants could be worse than no
mitigation at small budgets. Future work includes the \texttt{full} profile, correlated readout noise with M3~\cite{nation2021},
exponential and optimised-node extrapolation~\cite{krebsbach2022}, PEC and CDR, coupling-map-aware compilation, and hardware validation.

\paragraph{Data and code availability.} All code, configurations, raw data manifests and figure scripts are in the QEM-Bench repository
(MIT license). Figures 5, 7 and 10 (overhead, family and CI plots) are in \texttt{paper/figures/} and the supplement. Figures and tables are regenerated from saved data by \texttt{scripts/reproduce\_all.sh}.

\bibliographystyle{unsrtnat}
\bibliography{references}
\end{document}
